\chapter{Conclusiones}
\label{chap:conclusions}
\lettrine[]{A} continuación, se detallan los aprendizajes técnicos y arquitectónicos más relevantes, la integración de los conocimientos del grado en el desarrollo del proyecto, y las oportunidades de mejora y expansión identificadas tras su implementación.

\section{Lecciones aprendidas durante implementación}
\label{sec:lecciones_aprendidas}

\subsection{Decisiones arquitectónicas importantes}
\label{subsec:decisiones_arquitectonicas}

\subsubsection{Refactorización temprana}

La refactorización temprana basada en principios SOLID fue crucial:

\textbf{Beneficios obtenidos:}
\begin{itemize}
    \item \textbf{Improved maintainability}: Código más fácil de mantener y extender
    \item \textbf{Better testing}: Componentes aislados facilitaron testing
    \item \textbf{Reduced coupling}: Menor acoplamiento entre componentes
    \item \textbf{Enhanced debugging}: Debugging más eficiente con responsabilidades claras
\end{itemize}

\subsubsection{APIs como herramienta de desarrollo}

El desarrollo de APIs fue fundamental para el proceso:

\textbf{Valor aportado:}
\begin{itemize}
    \item \textbf{Internal testing}: APIs facilitaron testing exhaustivo
    \item \textbf{Error detection}: Revelaron errores internos no detectados
    \item \textbf{Architecture validation}: Validaron la arquitectura del sistema
    \item \textbf{External integration}: Facilitaron integración con herramientas externas
\end{itemize}

\subsection{Desafíos técnicos superados}
\label{subsec:desafios_superados}

\subsubsection{Threading y Qt}

La integración de threading con Qt presentó desafíos:

\textbf{Soluciones implementadas:}
\begin{itemize}
    \item \textbf{Qt Signals/Slots}: Comunicación thread-safe entre hilos
    \item \textbf{Worker threads}: Procesamiento en hilos separados
    \item \textbf{Event-driven architecture}: Arquitectura basada en eventos para PS4
    \item \textbf{Adapter pattern}: Adaptadores para aislar threading complexity
\end{itemize}

\subsubsection{Performance optimization}

Los datos de alta frecuencia requirieron optimizaciones específicas:

\textbf{Optimizaciones implementadas:}
\begin{itemize}
    \item \textbf{Worker pools}: Distribución de carga de procesamiento
    \item \textbf{Batch processing}: Procesamiento por lotes para eficiencia
    \item \textbf{Adaptive throttling}: Throttling inteligente basado en carga
    \item \textbf{Memory optimization}: Gestión eficiente de buffers y memoria
\end{itemize}

\subsection{Relación con competencias del grado}
\label{subsec:relacion_competencias}

\subsubsection{Aplicación de conocimientos fundamentales}

El desarrollo del framework ha integrado conocimientos de múltiples asignaturas del grado:

\textbf{Fundamentos de programación y arquitectura:}
\begin{itemize}
    \item \textbf{Diseño de Software}: Aplicación de paradigmas de programación orientada a objetos
    \item \textbf{Concurrencia y Paralelismo}: Breve pincelada a patrones orientados a eventos
    \item \textbf{Algoritmos + Gestión de interfaces}: Optimización de algoritmos de procesamiento de señales
\end{itemize}

\textbf{Desarrollo software y diseño:}
\begin{itemize}
    \item \textbf{Deseño Software}: Aplicación extensiva de principios SOLID y patrones de diseño
    \item \textbf{Bases de Datos}: Diseño de estructuras de datos para almacenamiento eficiente
    \item \textbf{Interfaces Persoa Máquina}: Desarrollo de interfaces intuitivas multi-modal
\end{itemize}

\textbf{Sistemas y redes:}
\begin{itemize}
    \item \textbf{Sistemas Operativos}: Gestión de procesos, hilos y sincronización
    \item \textbf{Concorrencia e Paralelismo}: Thread-safety y procesamiento multiproceso
    \item \textbf{Redes}: Implementación de comunicación WebSocket y protocols de red
    \item \textbf{Internet e Sistemas Distribuídos}: Arquitectura REST
\end{itemize}

\subsubsection{Competencias de mención en computación}

Las asignaturas de especialización han sido particularmente relevantes:

\textbf{Inteligencia artificial y procesamiento:}
\begin{itemize}
    \item \textbf{Sistemas Inteligentes + Aprendizaxe Automática}: Base para métricas
    \item \textbf{Visión Artificial}: Aunque no directamente aplicada a señales, se dieron multiples formas de procesar imágenes que se pueden extrapolar a señales
    \item \textbf{Procesamento de Linguaxes}: Desarrollo de CLI (parsing)
\end{itemize}

\subsection{Competencias transversales desarrolladas}
\label{subsec:competencias_transversales}

\subsubsection{Investigación y aprendizaje autónomo}

\textbf{Competencias desarrolladas:}
\begin{itemize}
    \item \textbf{Investigación bibliográfica}: Análisis del estado del arte en BCI
    \item \textbf{Aprendizaje de tecnologías}: Familiarización con nuevas librerías (BrainFlow, MNE-Python)
    \item \textbf{Resolución de problemas}
    \item \textbf{Pensamiento crítico}: Evaluación y selección de alternativas tecnológicas
\end{itemize}

\section{Líneas de trabajo futuro}
\label{sec:lineas_futuro}

\subsection{Mejoras técnicas identificadas}
\label{subsec:mejoras_tecnicas}

\subsubsection{Modo headless completo}

Una limitación importante identificada es la dependencia de Qt para el funcionamiento headless:

\textbf{Problemas actuales:}
\begin{itemize}
    \item \textbf{Dependencias Qt}: CLI y APIs requieren componentes Qt inicializados
    \item \textbf{Limitaciones de deployment}: Imposibilita uso en servidores sin display
    \item \textbf{Overhead de recursos}: Inicialización innecesaria de componentes gráficos
\end{itemize}

\textbf{Soluciones propuestas:}
\begin{itemize}
    \item \textbf{Refactorización de servicios}: Separar lógica de negocio de componentes Qt
    \item \textbf{Alternative backends}: Implementación de backends alternativos sin Qt
\end{itemize}

\subsubsection{Optimizaciones de rendimiento}

\textbf{Áreas de mejora identificadas:}
\begin{itemize}
    \item \textbf{Tree.js integration}: Migración de Chart.js a Tree.js para mejor rendimiento
    \item \textbf{CUDA acceleration}: Aceleración GPU para operaciones ICA y wavelet
    \item \textbf{Compression algorithms}: Compresión en tiempo real para reducir overhead de red
\end{itemize}